% Options for packages loaded elsewhere
\PassOptionsToPackage{unicode}{hyperref}
\PassOptionsToPackage{hyphens}{url}
\documentclass[
]{article}
\usepackage{xcolor}
\usepackage[margin=1in]{geometry}
\usepackage{amsmath,amssymb}
\setcounter{secnumdepth}{-\maxdimen} % remove section numbering
\usepackage{iftex}
\ifPDFTeX
  \usepackage[T1]{fontenc}
  \usepackage[utf8]{inputenc}
  \usepackage{textcomp} % provide euro and other symbols
\else % if luatex or xetex
  \usepackage{unicode-math} % this also loads fontspec
  \defaultfontfeatures{Scale=MatchLowercase}
  \defaultfontfeatures[\rmfamily]{Ligatures=TeX,Scale=1}
\fi
\usepackage{lmodern}
\ifPDFTeX\else
  % xetex/luatex font selection
\fi
% Use upquote if available, for straight quotes in verbatim environments
\IfFileExists{upquote.sty}{\usepackage{upquote}}{}
\IfFileExists{microtype.sty}{% use microtype if available
  \usepackage[]{microtype}
  \UseMicrotypeSet[protrusion]{basicmath} % disable protrusion for tt fonts
}{}
\makeatletter
\@ifundefined{KOMAClassName}{% if non-KOMA class
  \IfFileExists{parskip.sty}{%
    \usepackage{parskip}
  }{% else
    \setlength{\parindent}{0pt}
    \setlength{\parskip}{6pt plus 2pt minus 1pt}}
}{% if KOMA class
  \KOMAoptions{parskip=half}}
\makeatother
\usepackage{longtable,booktabs,array}
\usepackage{calc} % for calculating minipage widths
% Correct order of tables after \paragraph or \subparagraph
\usepackage{etoolbox}
\makeatletter
\patchcmd\longtable{\par}{\if@noskipsec\mbox{}\fi\par}{}{}
\makeatother
% Allow footnotes in longtable head/foot
\IfFileExists{footnotehyper.sty}{\usepackage{footnotehyper}}{\usepackage{footnote}}
\makesavenoteenv{longtable}
\usepackage{graphicx}
\makeatletter
\newsavebox\pandoc@box
\newcommand*\pandocbounded[1]{% scales image to fit in text height/width
  \sbox\pandoc@box{#1}%
  \Gscale@div\@tempa{\textheight}{\dimexpr\ht\pandoc@box+\dp\pandoc@box\relax}%
  \Gscale@div\@tempb{\linewidth}{\wd\pandoc@box}%
  \ifdim\@tempb\p@<\@tempa\p@\let\@tempa\@tempb\fi% select the smaller of both
  \ifdim\@tempa\p@<\p@\scalebox{\@tempa}{\usebox\pandoc@box}%
  \else\usebox{\pandoc@box}%
  \fi%
}
% Set default figure placement to htbp
\def\fps@figure{htbp}
\makeatother
\setlength{\emergencystretch}{3em} % prevent overfull lines
\providecommand{\tightlist}{%
  \setlength{\itemsep}{0pt}\setlength{\parskip}{0pt}}
\usepackage{bookmark}
\IfFileExists{xurl.sty}{\usepackage{xurl}}{} % add URL line breaks if available
\urlstyle{same}
\hypersetup{
  hidelinks,
  pdfcreator={LaTeX via pandoc}}

\author{}
\date{\vspace{-2.5em}}

\begin{document}

\section{📚 Collection Jeux de Rôle - Résolution de
Problèmes}\label{collection-jeux-de-ruxf4le---ruxe9solution-de-probluxe8mes}

\textbf{6 jeux de rôle complets pour la formation aux méthodes de
résolution de problèmes industriels}

\begin{center}\rule{0.5\linewidth}{0.5pt}\end{center}

\subsection{📖 COMPRENDRE
L'ORGANISATION}\label{comprendre-lorganisation}

Cette collection contient \textbf{2 types de fichiers} :

\subsubsection{1. Documentation Générale (À LA
RACINE)}\label{documentation-guxe9nuxe9rale-uxe0-la-racine}

\begin{longtable}[]{@{}
  >{\raggedright\arraybackslash}p{(\linewidth - 4\tabcolsep) * \real{0.2195}}
  >{\raggedright\arraybackslash}p{(\linewidth - 4\tabcolsep) * \real{0.3171}}
  >{\raggedright\arraybackslash}p{(\linewidth - 4\tabcolsep) * \real{0.4634}}@{}}
\toprule\noalign{}
\begin{minipage}[b]{\linewidth}\raggedright
Fichier
\end{minipage} & \begin{minipage}[b]{\linewidth}\raggedright
Description
\end{minipage} & \begin{minipage}[b]{\linewidth}\raggedright
Quand l'utiliser ?
\end{minipage} \\
\midrule\noalign{}
\endhead
\bottomrule\noalign{}
\endlastfoot
\textbf{\href{./GUIDE_NAVIGATION.md}{\texttt{GUIDE\_NAVIGATION.md}}} 🔑
& \textbf{COMMENCEZ ICI !} Guide pour comprendre l'organisation &
Première lecture \\
\href{./README.md}{\texttt{README.md}} & Vue d'ensemble de la collection
& Vue d'ensemble rapide \\
\href{./Jeux_de_Role_Resolution_Problemes.md}{\texttt{Jeux\_de\_Role\_Resolution\_Problemes.md}}
& \textbf{Descriptions des 6 scénarios} en 1 fichier & Découvrir les 6
jeux \\
\href{./Guide_Formateur_Jeux_de_Role.md}{\texttt{Guide\_Formateur\_Jeux\_de\_Role.md}}
& Méthodologie d'animation & Apprendre à animer \\
\href{./Index_Collection_Jeux_de_Role.md}{\texttt{Index\_Collection\_Jeux\_de\_Role.md}}
& Catalogue détaillé & Comparer et choisir \\
\end{longtable}

\subsubsection{2. Kits de Jeu (DANS LES
DOSSIERS)}\label{kits-de-jeu-dans-les-dossiers}

Chaque jeu de rôle a son propre dossier avec tout le matériel nécessaire
:

\begin{longtable}[]{@{}
  >{\raggedright\arraybackslash}p{(\linewidth - 4\tabcolsep) * \real{0.2727}}
  >{\raggedright\arraybackslash}p{(\linewidth - 4\tabcolsep) * \real{0.2727}}
  >{\raggedright\arraybackslash}p{(\linewidth - 4\tabcolsep) * \real{0.4545}}@{}}
\toprule\noalign{}
\begin{minipage}[b]{\linewidth}\raggedright
Dossier
\end{minipage} & \begin{minipage}[b]{\linewidth}\raggedright
Contenu
\end{minipage} & \begin{minipage}[b]{\linewidth}\raggedright
Kit complet ?
\end{minipage} \\
\midrule\noalign{}
\endhead
\bottomrule\noalign{}
\endlastfoot
\href{./JdR1_QRQC/}{\texttt{JdR1\_QRQC/}} & QRQC de survie & ✅ OUI \\
\href{./JdR2_Methode_8D/}{\texttt{JdR2\_Methode\_8D/}} & Méthode 8D en
crise & ✅ OUI \\
\href{./JdR3_5_Pourquoi/}{\texttt{JdR3\_5\_Pourquoi/}} & 5 Pourquoi du
mystère & ✅ OUI \\
\href{./JdR4_Pareto/}{\texttt{JdR4\_Pareto/}} & Pareto sous pression &
✅ OUI \\
\href{./JdR5_Poka_Yoke/}{\texttt{JdR5\_Poka\_Yoke/}} & Poka-Yoke
salvateur & ✅ OUI \\
\href{./JdR6_Kaizen/}{\texttt{JdR6\_Kaizen/}} & Kaizen Blitz & ✅ OUI \\
\end{longtable}

\begin{center}\rule{0.5\linewidth}{0.5pt}\end{center}

\subsection{🚀 DÉMARRAGE RAPIDE}\label{duxe9marrage-rapide}

\subsubsection{Option A : Je découvre la
collection}\label{option-a-je-duxe9couvre-la-collection}

\begin{enumerate}
\def\labelenumi{\arabic{enumi}.}
\tightlist
\item
  Lisez \href{./GUIDE_NAVIGATION.md}{\texttt{GUIDE\_NAVIGATION.md}} ←
  \textbf{Commencez ICI}
\item
  Parcourez
  \href{./Jeux_de_Role_Resolution_Problemes.md}{\texttt{Jeux\_de\_Role\_Resolution\_Problemes.md}}
  pour voir les 6 scénarios
\item
  Consultez
  \href{./Guide_Formateur_Jeux_de_Role.md}{\texttt{Guide\_Formateur\_Jeux\_de\_Role.md}}
  pour la méthodologie
\end{enumerate}

\subsubsection{Option B : Je veux animer un jeu spécifique (ex:
QRQC)}\label{option-b-je-veux-animer-un-jeu-spuxe9cifique-ex-qrqc}

\begin{enumerate}
\def\labelenumi{\arabic{enumi}.}
\tightlist
\item
  Ouvrez le dossier \href{./JdR1_QRQC/}{\texttt{JdR1\_QRQC/}}
\item
  Lisez \texttt{README.md} (résumé du scénario)
\item
  Imprimez \texttt{Kit\_Indices\_JDR1\_QRQC.md} (tous les indices +
  solution)
\item
  Animez !
\end{enumerate}

\begin{center}\rule{0.5\linewidth}{0.5pt}\end{center}

\subsection{📞 Contact}\label{contact}

\textbf{Auteur} : S. Jaubert\\
\textbf{Organisation} : Pôle Formation UIMM - Centre-Val de Loire\\
\textbf{Date} : Janvier 2026

\begin{center}\rule{0.5\linewidth}{0.5pt}\end{center}

\emph{Cette collection est prête à l'emploi pour la formation !} 🎭✨

\end{document}
